%==============================================================================
% Sjabloon poster bachproef
%==============================================================================
% Gebaseerd op document class `a0poster' door Gerlinde Kettl en Matthias Weiser
% Aangepast voor gebruik aan HOGENT door Jens Buysse en Bert Van Vreckem

\documentclass[a0,portrait]{hogent-poster}
\usepackage[most]{tcolorbox}
\usepackage{tikz} 
\usetikzlibrary{positioning,fit,arrows.meta}

% Info over de opleiding
\course{Bachelorproef}
\studyprogramme{toegepaste informatica}
\academicyear{2025-2026}
\institution{Hogeschool Gent, Valentin Vaerwyckweg 1, 9000 Gent}

% Info over de bachelorproef
\title{Burn-outpreventie bij Vlaamse twintigers: een vergelijkende studie van bestaande apps en vertaling naar ontwerpprincipes voor een stressreducerende applicatie}
\subtitle{}
\author{Britt Emanuel}
\email{britt.emanuel@student.hogent.be}
\supervisor{Chloé De Leenheer}
\cosupervisor{Katrijn Van den Broeck (Akompanje)}

% Indien ingevuld, wordt deze informatie toegevoegd aan het einde van de
% abstract. Zet in commentaar als je dit niet wilt.
\specialisation{Full stack developer}
\keywords{M-health, Cognitieve belasting, Stressreducerend ontwerp }
% \projectrepo{https://github.com/user/repo}

\begin{document}

\maketitle

\begin{abstract}
Deze bachelorproef onderzoekt hoe digitale toepassingen kunnen bijdragen aan stress- en burn-outpreventie bij Vlaamse twintigers. Via een literatuurstudie, een expertinterview en een requirementsanalyse werden concrete ontwerpprincipes opgesteld voor stressreducerende ontwerp. Vervolgens werd een vergelijkende studie van bestaande mentale gezondheidsapps uitgevoerd om deze ontwerpprincipes te toetsen en aanvullende inzichten te verzamelen die konden worden meegenomen in de verdere analyse en ontwikkeling van de Proof-of-Concept. Op basis van deze inzichten werd een cross-platform Proof-of-Concept ontwikkeld. Het onderzoek toont aan dat naast de inhoud van een mentale gezondheidsapp ook aandacht voor psychologische basisbehoeften, personalisatie, eenvoud, beperkte cognitieve belasting en controle over digitale prikkels belangrijke factoren zijn voor een stressreducerende gebruikerservaring.
\end{abstract}

\begin{multicols}{2} % This is how many columns your poster will be broken into, a portrait poster is generally split into 2 columns

\section{Introductie}

\vspace{1em}

Vandaag de dag blijft het aantal jongvolwassenen dat te maken krijgt met stress en burn-out alsmaar toenemen. Volgens cijfers van het RIZIV steeg het aantal personen jonger dan 30 jaar dat langdurig arbeidsongeschikt is door burn-out van 825 in 2018 naar 1.707 in 2023, een toename van 107\%.

\vspace{1em}

Om hierop in te spelen en ondersteuning toegankelijker te maken, ontwikkelden steeds meer developers digitale toepassingen voor mentaal welzijn. Dit leidde tot een sterk groeiend aanbod: in 2019 waren er al meer dan 1.400 mentale gezondheidsapps beschikbaar. Voor een generatie die dagelijks met digitale technologie omgaat, lijken deze toepassingen een laagdrempelige manier om ondersteuning te bieden bij stress en burn-outpreventie. Toch blijft de burn-outproblematiek bij jongvolwassenen toenemen, wat de vraag oproept in welke mate bestaande apps daadwerkelijk aansluiten bij hun noden.

\vspace{1em}

\begin{tcolorbox}[
    width=\columnwidth,
    colback=green!10,
    colframe=green,
    boxrule=1pt,
    arc=6pt,
    outer arc=6pt,
    left=20pt,
    right=20pt,
    top=15pt,
    bottom=15pt,
    enhanced
]
    \medium
    \textbf{Onderzoeksvraag}: 
    Hoe kunnen apps voor stress- en burn-outpreventie het mentale welzijn van Vlaamse twintigers verbeteren, door middel van inzichten uit een analyse van bestaande apps te vertalen naar concrete ontwerpprincipes en functionaliteiten voor een proof-of-concept app?
\end{tcolorbox}

\vspace{1em}

Binnen het \textbf{probleemdomein} werd onderzocht welke noden en risicofactoren Vlaamse twintigers ervaren, waar bestaande apps tekortschieten en welke functionaliteiten en ontwerpprincipes worden toegepast.

\vspace{1em}

Binnen het \textbf{oplossingsdomein} werd gekeken naar essentiële requirements, UI/UX-elementen, cognitieve belasting, notificaties en gamification, en hoe deze kunnen worden vertaald naar een stressreducerend ontwerp.

\vspace{1em}

\section{Methodologie}

\vspace{1em}

\begin{center}
    \captionsetup{type=figure}
    \includegraphics[width=0.9\linewidth]{methodologie_schema}
    \captionof{figure}{Overzicht van de verschillende methodologische fases en hun onderlinge samenhang.}
\end{center}

\columnbreak


\section{Resultaten}

\vspace{1em}

Uit de analyse blijkt dat bestaande apps sterk inzetten op het aanbieden van tools, maar daarbij onvoldoende rekening houden met de mentale rust van de gebruiker. De belangrijkste aandachtspunten:

\vspace{1em}

\begin{center}
    \noindent
    \begin{minipage}[t]{0.235\columnwidth}
        \begin{tcolorbox}[
            width=\linewidth,
            height=5.2cm,
            valign=center,
            colback=hogent-blue!10,
            colframe=hogent-blue,
            boxrule=1pt,
            arc=6pt,
            boxsep=8pt
            ]
            \centering
            {\medium\textbf{1}}\\[0.4em]
            {\medium hoge\\cognitieve\\belasting}
        \end{tcolorbox}
    \end{minipage}%
    \hfill
    \begin{minipage}[t]{0.235\columnwidth}
        \begin{tcolorbox}[
            width=\linewidth,
            height=5.2cm,
            valign=center,
            colback=gray!10,
            colframe=gray!60,
            boxrule=1pt,
            arc=6pt,
            boxsep=8pt
            ]
            \centering
            {\medium\textbf{2}}\\[0.4em]
            {\medium beperkte\\personalisatie\\mogelijkheden}
        \end{tcolorbox}
    \end{minipage}%
    \hfill
    \begin{minipage}[t]{0.235\columnwidth}
        \begin{tcolorbox}[
            width=\linewidth,
            height=5.2cm,
            valign=center,
            colback=orange!10,
            colframe=orange,
            boxrule=1pt,
            arc=6pt,
            boxsep=8pt
            ]
            \centering
            {\medium\textbf{3}}\\[0.4em]
            {\medium prestatie\\gerichte\\motivatie}
        \end{tcolorbox}
    \end{minipage}%
    \hfill
    \begin{minipage}[t]{0.235\columnwidth}
        \begin{tcolorbox}[
            width=\linewidth,
            height=5.2cm,
            valign=center,
            colback=hogent-green!10,
            colframe=hogent-green!70!black,
            boxrule=1pt,
            arc=6pt,
            boxsep=8pt
            ]
            \centering
            {\medium\textbf{4}}\\[0.4em]
            {\medium beperkte\\notificatie\\controle}
        \end{tcolorbox}
    \end{minipage}
\end{center}

\vspace{1em}

\textbf{19 ontwerpprincipes} werden afgeleid uit de literatuurstudie, het expertinterview en de analyse van bestaande toepassingen. Ze werden gegroepeerd rond \textbf{development, design, interaction} en \textbf{content}. Ze vertalen psychologische en UX-inzichten naar concrete richtlijnen voor stressreductie. %Bijvoorbeeld:

%\begin{tcolorbox}[
%    width=\columnwidth,
%    colback=orange!10,
%    colframe=orange,
%    boxrule=1pt,
%    arc=6pt,
%    outer arc=6pt,
%    left=15pt,
%    right=15pt,
%    top=10pt,
%    bottom=10pt,
%    enhanced
%    ]
%    \begin{itemize}
%        \item \textbf{Ondersteun autonomie}: \textit{"Kan de gebruiker zelf keuzes maken en controle uitoefenen over zijn gebruik van de applicatie?"}
%        \item \textbf{Ondersteun competentiegevoel}: \textit{"Ervaart de gebruiker voldoende succeservaringen en begrijpt hij het effect van zijn acties?"}
%        \item \textbf{Stimuleer authentieke verbondenheid}: \textit{"Kan de gebruiker op een veilige en ondersteunende manier ervaringen delen zonder sociale druk?"}
%    \end{itemize}
%\end{tcolorbox}

\vspace{1em}

\begin{center}
    \captionsetup{type=figure}
    \includegraphics[width=1.0\linewidth]{PoC_schema}
    \captionof{figure}{Vergelijking van het aantal aandachtspunten, toepassingen en aanwezige functionele requirements (FR's) voor en na toepassing van de ontwerpprincipes.}
\end{center}

\vspace{1em}

\section{Conclusie \& toekomstig onderzoek}

\vspace{1em}

Het onderzoek toont aan dat een stressreducerende m-health app niet alleen draait om de interventies, maar ook om \textbf{hoe deze worden aangeboden}. De analyse resulteerde in \textbf{19 ontwerpprincipes} rond onder andere personalisatie, lage cognitieve belasting en ondersteunende communicatie. Samen vormen ze een \textbf{praktisch kader} voor het ontwerpen en beoordelen van stressreducerende apps. 

\begin{tcolorbox}[
    width=\columnwidth,
    colback=hogent-blue!10,
    colframe=hogent-blue,
    boxrule=1pt,
    arc=6pt,
    outer arc=6pt,
    left=20pt,
    right=20pt,
    top=15pt,
    bottom=15pt,
    enhanced
    ]
    \medium
    \textbf{Vervolgonderzoek} moet nagaan of de ontwerpprincipes duidelijk, eenduidig en toepasbaar zijn voor andere ontwikkelaars. Alsook kan bekeken worden hoe het principe \textbf{toegankelijkheid} concreter gedefinieerd kan worden. Daarnaast kan een verbeterde app langdurig met Vlaamse twintigers worden getest op de bijdrage aan het verminderen van chronische overspanning.
\end{tcolorbox}

\vspace{1em}

\section{Referenties}

De bp, samen met de bijhorende PoC, is beschikbaar via de volgende GitHub-repository:
\url{https://github.com/BrittEmanuel2001/Bachelorproef-25-26-BE}



\end{multicols}
\end{document}